\usepackage{amsmath, amssymb, amsfonts}

\usepackage{amsthm}
\newtheorem{definition}{Định nghĩa}[section]
\newtheorem{theorem}{Định lý}[section]

\usepackage{bm}
\usepackage{mathrsfs}
\usepackage{graphicx}
\usepackage{setspace}
\usepackage{caption}
\usepackage{subcaption}
\usepackage{float}
\usepackage{titlesec}
\usepackage{xcolor}
\usepackage{xstring}
\usepackage{listings}
\usepackage{hyperref}
\usepackage{tikz}
\usepackage{booktabs}
\usepackage{adjustbox}


% appendix
\usepackage{nomencl} 
\usepackage{etoolbox} % Định nghĩa các nhóm trong Danh mục

\usepackage{tikz}
% --- THÊM DÒNG NÀY ĐỂ SỬA LỖI ---
% --------------------------------

\usepackage{booktabs, multirow, multicol}
\usepackage{diagbox} % Dùng để kẻ ô chéo trong bảng vận tải (góc Tây - Bắc)

\usepackage{emptypage} % Xóa số trang ở các trang trắng tự động
\usepackage{enumitem} % Dùng để tùy chỉnh khoảng cách các danh sách

\usetikzlibrary{
  arrows.meta, 
  positioning, 
  calc, 
  shapes.geometric, 
  shadows,
  matrix,
  positioning,
  fit,
  arrows.meta
  }


% =========================
% GÓI LỆNH VÀ LỆNH TỰ TẠO CHO BẢNG ĐỐI NGẪU
% =========================
\usepackage{cancel} % Dùng để gạch chéo các dấu không thỏa mãn

% 1. Lệnh chèn ký hiệu góc trên bên trái (Top-Left) và giá trị vào ô
\newcommand{\TL}[2]{%
  \tikz[baseline=(val.base)]{
    % 1. Ký hiệu góc trên
    \node[anchor=south west, inner sep=0pt, font=\scriptsize, text=gray] (label) {$#1$};
    
    % 2. Giá trị chính nằm ngay dưới ký hiệu (yshift điều chỉnh khoảng cách dòng)
    \node[anchor=north west, inner sep=0pt, font=\normalsize\bfseries, text=black] (val) at ([yshift=-1.5mm]label.south west) {#2};
    
    % 3. Tạo một khung bao tàng hình (bounding box) có khoảng đệm 4pt để ép LaTeX nới rộng hàng, không chạm viền bảng
    \path ([shift={(-4pt, 4pt)}]label.north west) rectangle ([shift={(4pt, -4pt)}]val.south east);
  }%
}

% 2. Lệnh khoanh tròn dấu (dùng fbox để đóng khung)
\newcommand{\khoanh}[1]{\setlength{\fboxsep}{3pt}\fbox{#1}}
% =========================

% Định nghĩa lệnh khoanh tròn màu sắc
\newcommand{\kh}[2]{\tikz[baseline=(char.base)] \node[shape=circle, draw=#1, fill=#1!10, text=black, inner sep=1.5pt, font=\bfseries] (char) {#2};}
% Định nghĩa lệnh khoanh tròn nét đôi (phù hợp in cả màu và trắng đen)
\newcommand{\khdouble}[2]{\tikz[baseline=(char.base)] \node[shape=circle, draw=#1, double, double distance=1pt, text=black, inner sep=1.5pt, font=\bfseries] (char) {#2};}
% Dash
\newcommand{\khdash}[2]{\tikz[baseline=(char.base)] \node[shape=circle, draw=#1, dashed, text=black, inner sep=1.5pt, font=\bfseries] (char) {#2};}

% =========================================================
% BỘ LỆNH ĐÓNG KHUNG TIKZ DÀNH CHO QUY HOẠCH TUYẾN TÍNH
% Cần đảm bảo đã gọi: \usepackage{tikz} 
\usetikzlibrary{shapes.geometric} % (nếu dùng hình oval)
% =========================================================
% 0. Rectangle
\newcommand{\kv}[2]{\tikz[baseline=(char.base)] \node[shape=rectangle, draw=#1, fill=#1!10, text=black, inner sep=1.5pt, font=\bfseries] (char) {#2};}
% 1. Khung nét đôi
\newcommand{\kvd}[2]{\tikz[baseline=(char.base)] \node[shape=rectangle, draw=#1, double, double distance=1pt, text=black, inner sep=2pt, font=\bfseries] (char) {#2};}
% 2. Khung nét đứt (Dashed) - Thường dùng cho hệ số ngoài cơ sở
\newcommand{\kdash}[2]{\tikz[baseline=(char.base)] \node[shape=rectangle, draw=#1, dashed, text=black, inner sep=2pt, font=\bfseries] (char) {#2};}
% 3. Khung nét chấm bi (Dotted) - Thường dùng cho hệ số cơ sở
\newcommand{\kdot}[2]{\tikz[baseline=(char.base)] \node[shape=rectangle, draw=#1, dotted, thick, text=black, inner sep=2pt, font=\bfseries] (char) {#2};}
% 4. Khung nét đậm (Thick) - Dùng để nhấn mạnh phần tử trụ (Pivot)
\newcommand{\kthick}[2]{\tikz[baseline=(char.base)] \node[shape=rectangle, draw=#1, very thick, text=black, inner sep=2pt, font=\bfseries] (char) {#2};}
% 5. Khung bo tròn góc (Rounded) - Nhìn mềm mại hơn khung vuông sắc cạnh
\newcommand{\kround}[2]{\tikz[baseline=(char.base)] \node[shape=rectangle, draw=#1, rounded corners=3pt, text=black, inner sep=2pt, font=\bfseries] (char) {#2};}
% 6. Khung có màu nền nhạt (Fill) - Dùng làm nổi bật hẳn 1 ô (in xám trắng vẫn thấy nền)
\newcommand{\kfill}[2]{\tikz[baseline=(char.base)] \node[shape=rectangle, draw=#1, fill=#1!15, text=black, inner sep=2pt, font=\bfseries] (char) {#2};}
% 7. Khung hình Oval - Dùng cho cột Phương án (PA) như sơ đồ trong giáo trình
\newcommand{\koval}[2]{\tikz[baseline=(char.base)] \node[shape=ellipse, draw=#1, text=black, inner sep=1.5pt, font=\bfseries] (char) {#2};}


\usepackage[most]{tcolorbox}
% \tcbuselibrary{theorems, ams}

% Định nghĩa một môi trường mới cho phương trình có khung
\newtcolorbox{formula}[1][]{
    ams equation,       % Tích hợp môi trường equation của amsmath
    enhanced,           % Kích hoạt các tính năng đồ họa nâng cao
    colback=white,      % Màu nền trắng (phù hợp in ấn học thuật)
    colframe=black,     % Màu viền đen
    boxrule=0.5pt,      % Độ dày viền tiêu chuẩn
    sharp corners,      % Góc vuông để giữ vẻ chuyên nghiệp
    fonttitle=\bfseries,
    title=#1,           % Tham số tùy chọn để đặt tên cho công thức
    before skip=10pt,   % Khoảng cách phía trên box
    after skip=10pt     % Khoảng cách phía dưới box
}

\newcommand{\LPsemanticheading}[3]{%
  \par\medskip
  \LPsidebarlabel{#1}%
  {\noindent\color{#2}\rule{\linewidth}{0.5pt}\par}
  {\noindent\bfseries\color{#2} #3\par}
  {\noindent\color{#2!70!black}\rule{0.38\linewidth}{0.4pt}\par\smallskip}
}

\NewDocumentEnvironment{exampleblock}{m}{%
  \LPsemanticheading{Ví dụ}{blue!60!black}{#1}%
}{\par\medskip}

\NewDocumentEnvironment{exerciseblock}{m}{%
  \LPsemanticheading{Bài tập}{orange!70!black}{#1}%
}{\par\medskip}

\NewDocumentEnvironment{questionblock}{m}{%
  \LPsemanticheading{Câu hỏi}{red!70!black}{#1}%
}{\par\medskip}

\NewDocumentEnvironment{remarkblock}{m}{%
  \LPsemanticheading{Ghi chú}{green!50!black}{#1}%
}{\par\medskip}

\NewDocumentEnvironment{solutionblock}{m}{%
  \LPsemanticheading{Lời giải}{black!70}{#1}%
}{\par\medskip}

\makenomenclature

% =========================
% FONT TIẾNG VIỆT
% =========================
\usepackage{iftex}

\ifPDFTeX
\usepackage[T5]{fontenc}
\usepackage[utf8]{inputenc}
\usepackage[vietnamese]{babel}
\else
\usepackage{fontspec}
\usepackage{polyglossia}
\setdefaultlanguage{vietnamese}
\defaultfontfeatures{Ligatures=TeX}
\IfFontExistsTF{Times New Roman}{
  \setmainfont{Times New Roman}
}{
  \setmainfont{TeX Gyre Termes}
}
\fi

\usepackage{geometry} % Chỉ gọi tên gói
\geometry{
  a4paper,
  top=2cm,
  bottom=2cm,
  inner=2.5cm,
  outer=0.7cm,
  marginparsep=0.9cm,
  marginparwidth=4.5cm,
  includemp,
  twoside
}

\raggedbottom

% ===== BỐ CỤC 3/7: CỘT GHI CHÚ + NỘI DUNG CHÍNH =====
\newif\ifLPshowsidebarlayout
\LPshowsidebarlayouttrue

\normalmarginpar
\setlength{\marginparpush}{10pt}
\newlength{\LPsidebaroffset}
\setlength{\LPsidebaroffset}{\dimexpr0.7cm+4.5cm+0.5\marginparsep\relax}

\definecolor{sidebarline}{gray}{0.75}
\definecolor{sidebartext}{gray}{0.30}

\newcommand{\LPsidebarlabel}[1]{%
  \marginpar{%
    \raggedright
    \sffamily
    \footnotesize
    \color{sidebartext}
    \rule{\marginparwidth}{0.4pt}\par
    \vspace{2pt}
    \textbf{\MakeUppercase{#1}}\par
    \vspace{2pt}
    \rule{0.62\marginparwidth}{0.4pt}%
  }%
}

\newcommand{\LPheadinglabel}[1]{%
  \begingroup
  \IfSubStr{#1}{Ví dụ}{\LPsidebarlabel{Ví dụ}}{%
    \IfSubStr{#1}{Bài tập}{\LPsidebarlabel{Bài tập}}{%
      \IfSubStr{#1}{Đề bài}{\LPsidebarlabel{Câu hỏi}}{%
        \IfSubStr{#1}{Đề thi mẫu}{\LPsidebarlabel{Đề thi}}{%
          \IfSubStr{#1}{ĐÁP ÁN}{\LPsidebarlabel{Đáp án}}{%
            \IfSubStr{#1}{Đáp án}{\LPsidebarlabel{Đáp án}}{}%
          }%
        }%
      }%
    }%
  }%
  \endgroup
}

\BeforeBeginEnvironment{definition}{\LPsidebarlabel{Định nghĩa}}
\BeforeBeginEnvironment{theorem}{\LPsidebarlabel{Định lý}}
\BeforeBeginEnvironment{formula}{\LPsidebarlabel{Công thức}}

\AddToHook{shipout/background}{%
  \ifLPshowsidebarlayout
    \begin{tikzpicture}[remember picture,overlay]
      \ifodd\value{page}
        \draw[line width=0.4pt, sidebarline]
          ([xshift=-\LPsidebaroffset,yshift=-2cm]current page.north east) --
          ([xshift=-\LPsidebaroffset,yshift=2cm]current page.south east);
      \else
        \draw[line width=0.4pt, sidebarline]
          ([xshift=\LPsidebaroffset,yshift=-2cm]current page.north west) --
          ([xshift=\LPsidebaroffset,yshift=2cm]current page.south west);
      \fi
    \end{tikzpicture}%
  \fi
}

% ===== TIÊU ĐỀ CÁC CẤP MỤC =====
\setcounter{secnumdepth}{3}
\setcounter{tocdepth}{3}

\renewcommand{\chaptername}{Chương}
\renewcommand{\thesubsubsection}{%
  \ifnum\value{subsection}=0
    \thesection.\arabic{subsubsection}%
  \else
    \thesubsection.\arabic{subsubsection}%
  \fi
}

\titleformat{\chapter}[display]
  {\normalfont\bfseries\centering}
  {\Large \chaptername~\thechapter}
  {0.5ex}
  {\Huge}

\titleformat{\section}
  {\normalfont\Large\bfseries}
  {Mục~\thesection.}
  {0.75em}
  {}

\titleformat{\subsection}
  {\normalfont\large\bfseries}
  {Tiểu mục~\thesubsection.}
  {0.75em}
  {}

\titleformat{\subsubsection}
  {\normalfont\normalsize\bfseries}
  {Mục nhỏ~\thesubsubsection.}
  {0.75em}
  {}

\titlespacing*{\chapter}{0pt}{0pt}{18pt}
\titlespacing*{\section}{0pt}{18pt}{8pt}
\titlespacing*{\subsection}{0pt}{14pt}{6pt}
\titlespacing*{\subsubsection}{0pt}{10pt}{4pt}

\let\LPsection\section
\RenewDocumentCommand{\section}{s o m}{%
  \LPheadinglabel{#3}%
  \IfBooleanTF{#1}
    {\LPsection*{#3}}
    {\IfNoValueTF{#2}{\LPsection{#3}}{\LPsection[#2]{#3}}}%
}

\let\LPsubsection\subsection
\RenewDocumentCommand{\subsection}{s o m}{%
  \LPheadinglabel{#3}%
  \IfBooleanTF{#1}
    {\LPsubsection*{#3}}
    {\IfNoValueTF{#2}{\LPsubsection{#3}}{\LPsubsection[#2]{#3}}}%
}

\let\LPsubsubsection\subsubsection
\RenewDocumentCommand{\subsubsection}{s o m}{%
  \LPheadinglabel{#3}%
  \IfBooleanTF{#1}
    {\LPsubsubsection*{#3}}
    {\IfNoValueTF{#2}{\LPsubsubsection{#3}}{\LPsubsubsection[#2]{#3}}}%
}

% ===== RENAME: Đổi tên danh mục =====
\renewcommand{\contentsname}{MỤC LỤC}
\renewcommand{\listfigurename}{DANH MỤC HÌNH ẢNH}
\renewcommand{\listtablename}{DANH MỤC BẢNG BIỂU}
\renewcommand{\nomname}{DANH MỤC KÝ HIỆU}

